Use o mesmo tipo de argumento energetico empregado no capitulo. Se um animal de tamanho caracteristico $L$ e geometricamente semelhante aos demais, entao:

\begin{itemize}
\item a massa corporal escala como $m\sim \rho L^3$;
\item a forca muscular maxima escala como area de secao, isto e, $F\sim \sigma L^2$, onde $\sigma$ e uma tensao muscular caracteristica.
\end{itemize}

Ao acelerar o corpo ao longo de uma passada de comprimento proporcional a $L$, o trabalho mecanico disponivel e da ordem
\[
W\sim FL\sim \sigma L^3.
\]
Esse trabalho alimenta a energia cinetica maxima,
\[
\frac{1}{2}mv_{\max}^2\sim \rho L^3 v_{\max}^2.
\]

Igualando as ordens de grandeza,
\[
\rho L^3 v_{\max}^2\sim \sigma L^3,
\]
e, portanto,
\[
v_{\max}^2\sim \frac{\sigma}{\rho},
\qquad
v_{\max}\sim \sqrt{\frac{\sigma}{\rho}}.
\]

O comprimento $L$ cancela completamente. Assim, em primeira aproximacao de semelhanca dinamica, a velocidade maxima de corrida e independente do tamanho do animal.

